% ====================================================================
% SECCIÓN MODULAR: PÁGINAS PRELIMINARES DE LA TESIS
% Requisito de la pauta: Incluir dedicatoria, agradecimientos y abstracts
% ====================================================================

% --- Nota: Dejar la dedicatoria lista al final, por ahora va este borrador ---
\thispagestyle{empty}
\vspace*{5cm}
\begin{flushright}
    \textit{Dedicado a mi familia, amigos y a quienes apoyaron y se sientan parte, \\ también a los muchos que faltarán.}
\end{flushright}

\newpage

% --- Sección de Agradecimientos ---
\section*{Agradecimientos}
% Nota: Texto de ejemplo.
Agradezco a la Universidad Católica del Maule y a la Facultad de Ciencias Básicas 

\newpage

% --- Resumen del proyecto ---
\section*{Resumen}
El hierro (Fe) constituye un micronutriente esencial para diversos procesos biogeoquímicos marinos y puede influir sobre la productividad primaria y la eficiencia de la bomba biológica de carbono. En este contexto, el aporte de hierro desde fuentes continentales y su posterior movilización y transporte constituyen procesos relevantes para comprender la disponibilidad de este elemento en los sistemas acuáticos. Sin embargo, la identificación de una fuente potencial de hierro requiere considerar no solo su composición, sino también las características del material y las condiciones ambientales que pueden favorecer su movilización y transporte.

El presente trabajo tiene como objetivo evaluar el potencial de fuentes geológicas de hierro y sus posibles rutas de transporte en un área de estudio acotada de la Región del Maule, mediante la integración de antecedentes geológicos, geoquímicos, geomorfológicos, topográficos, hidrológicos y climáticos. Para ello, se propone desarrollar una base de datos espacial que permita caracterizar potenciales fuentes de hierro y analizar su conectividad mediante posibles rutas de transporte fluvial y atmosférico, considerando diferentes condiciones de movilización y las incertidumbres asociadas a la información disponible.

Durante esta etapa se realiza una revisión bibliográfica orientada a establecer el marco conceptual del estudio, abordando la bomba de tejidos blandos, la función del hierro en la biogeoquímica oceánica, las características de las partículas minerales y los procesos asociados a su movilización y transporte. Asimismo, se contempla la recopilación y organización de información geoespacial y temática necesaria para la construcción de la base de datos y el posterior análisis espacial. La metodología propuesta permitirá desarrollar una aproximación exploratoria al sistema fuente--transporte--receptor, orientada a identificar zonas con condiciones potencialmente favorables para constituir fuentes de hierro y sus posibles conexiones con sistemas fluviales y ambientes costeros.

\vspace{1.5cm} % Para separar ordenadamente ambos textos

% --- Abstract en inglés (Cambiando el idioma con babel) ---
%\selectlanguage{english}
%\section*{Abstract}
%Texto aqui.
\selectlanguage{spanish} % volver al español aquí para no romper el formato